% Nouveau EGA V (cote 25): archive cote 25, site page 37 (PDF page 38).
% typescript. Draft compared: none.
% AI-transcribed by Claude Opus 5.5 (claude-opus-5-5, adaptive thinking, effort high), Message Batches API, batch ega25b (msgbatch_01A7n5YuXubUPNK1zsepi8RG), 2026-11-13; not yet checked by eye.
\scanpage{25/37}
\mspage{\hand{\unsure{39|33}}}

8.\var{\over{4}{6}}{6}{writing over}{ega25_25_037-v0001}. Nous étudions maintenant les points  y de $\underline{Y}$ qui sont non lisses
sur $Y_{\hand{\xi}}$ rel à $k(\hand{\xi})$. Nous nous  bornons au cas où  $f:X\hand{\to} P$  est \uline{non ramifié}
(pratiquement, ce sera une immersion), et où $X\hand{\to} S$ est lisse.  Nous ne
supposons pas nécessairement S le spectre d'un corps.  \var{\xout{Nous nous bornons}}{}{typewriter strike-out}{ega25_25_037-v0002}
\var{\xout{à la considération des points de $X_P$ qui sont dans l'ouvert U de}}{}{typewriter strike-out}{ega25_25_037-v0003} Comme
f  est non ramifié, l'homomorphisme canonique   $f^*(\hand{\underline{\Omega}}^1_{P/S})\hand{\to}\hand{\underline{\Omega}}^1_{X/S}$
est surjectif, et son noyau est un Module loc libre sur X, que nous
noterons  $\hand{\check{\nu}}_{X/P}$ ; lorsque f est une immersion, ce n'est autre que le
Module \uline{conormal}  $\underline{J}/\underline{J}^2$  défini par l'Idéal $\underline{J}$ de X dans P , et nous l'appe-%
lerons en tous cas Module conormal:
\[ \hand{(a)} \qquad 0\hand{\to}\hand{\check{\underline{\nu}}}_{X/P}\hand{\to} f^*(\hand{\underline{\Omega}}^1_{P/S})\hand{\to}\hand{\underline{\Omega}}^1_{X/S}\hand{\to} 0 \ . \]
\var{\margin{\hand{$|||$}}}{\placedmargin{\hand{$|||$}}}{marginal note placed}{ega25_25_037-v0004}
% NOTE: Three parallel pencil sidestrokes stand beside the canonical exact-sequence paragraph. Notons que nous avons aussi sur  P  une suite exacte canonique
(qui pouvait figurer à titre d'exemple au par.16 p.ex.)
\[ \hand{(b)} \qquad 0\hand{\to}\hand{\underline{\Omega}}^1_{P/S}(1)\hand{\to}\underline{E}_P\hand{\to}\underline{O}_P(1)\hand{\to} 0 \]
(i.e. $\hand{\underline{\Omega}}^1_{P/S}$ est can isomorphe au noyau de l'homomorphisme canonique
$\underline{E}_P(-1)\hand{\to}\underline{O}_P$  déduit de $\underline{E}_P\hand{\to}\underline{O}_P(1)$ ), d'où en appliquant $f^*$ :
\[ \hand{(b')} \qquad 0 \hand{-} f^*(\hand{\underline{\Omega}}^1_{P/S})(1) \hand{-} \underline{E}_X \hand{-} \underline{O}_X(1)\hand{\to} 0 \ \hand{\unsure{)|,}} \]
qui donne une description explicite de  $f^*(\hand{\underline{\Omega}}^1_{P/S})$  \var{\xout{\unsure{de X|du X}}}{}{typewriter strike-out}{ega25_25_037-v0005} sur X, et
permet  donc d'identifier $\hand{\check{\underline{\nu}}}_{X/P}(1)$  à un sous-Module loc facteur direct
de $\underline{E}_X$ , ou encore le \var{\xout{Module conormal}}{}{typewriter strike-out}{ega25_25_037-v0006} dual  $\hand{\underline{\nu}}_{X/P}(-1)$  à un Module
quotient de $\overset{\hand{\vee}}{\underline{E}}_X$ .  Par suite,   $P(\hand{\underline{\nu}}_{X/P}\var{\deleted{\unsure{*}}}{}{deletion}{ega25_25_037-v0007}(-1)\ )$  =  $P(\hand{\underline{\nu}}_{X/P})$  se plonge
canoniquement dans  $P(\overset{\hand{\vee}}{\underline{E}}_X) = X\times_S\overset{\hand{\vee}}{P} = X_{\overset{\hand{\vee}}{P}}$ comme un sous-fibré projectif sur
X, donc comme sous-préschéma fermé. Ce dernier est nécessairement contenu
dans $\underline{Y}$  (du fait que $\hand{\underline{\Omega}}_{X/P}(1)$  est contenu dans le noyau de  $\underline{E}_X\hand{\to}\underline{O}_X(1)$)
\var{\deleted{L'ensemble sous-jacent  à ce sous-préschéma n'est autre que l'ensemble des}}{}{deletion}{ega25_25_037-v0008}
\var{\deleted{\var{\xout{xxxxP}}{}{typewriter strike-out}{ega25_25_037-v0010} $y\in\overline{Y}$ \var{\xout{tel que}}{}{typewriter strike-out}{ega25_25_037-v0011} \uline{qui sont} \var{\inserted{\var{\deleted{\unsure{des zéros x singuliers pour}}}{}{deletion}{ega25_25_037-v0013}}}{\var{\deleted{\unsure{des zéros x singuliers pour}}}{}{deletion}{ega25_25_037-v0013}}{insertion}{ega25_25_037-v0012} \var{\deleted{"singuliers sur $Y_{\hand{\xi}}$ " au sens suivant:}}{}{deletion}{ega25_25_037-v0014}}}{}{deletion}{ega25_25_037-v0009}
\var{\deleted{\var{\margin{des zéros singuliers}}{\placedmargin{des zéros singuliers}}{marginal note placed}{ega25_25_037-v0016} \var{\deleted{un point est dit "non singulier" sur $Y_{\hand{\xi}}$ si $\hand{\phi_\xi}$ est \unsure{Dégulier|régulier}}}{}{deletion}{ega25_25_037-v0017} sur $X_{k(\hand{\xi})}$}}{}{deletion}{ega25_25_037-v0015}
\var{\deleted{en y et \var{\xout{\illegible{1}}}{}{typewriter strike-out}{ega25_25_037-v0019} $Y_{\hand{\xi}}$ est lisse sur $k(\hand{\xi})$ en y.}}{}{deletion}{ega25_25_037-v0018}

\mspage{\hand{37}}
